\section{What the Church Confesses}

A Catholic does not receive the teaching on hell as a private inference from scattered verses. He receives it as an article of the Catholic faith, confessed in creeds, defined by ecumenical councils and popes, prayed at the altar, and handed down without interruption from the apostolic age to the present day. Every solemn pronouncement of the Church on this subject has been called forth by a specific pressure: either to guard the eternity of the sentence against theories of temporary chastisement, or to guard the immediacy of retribution after death, or to guard both the justice of God and the true freedom of man.

Already in the early centuries, the Athanasian Creed---\textit{Quicumque vult}, received across the Latin West as a standard of orthodox faith---sets the resurrection and the twofold eternal destiny side by side in lapidary Latin: \textit{Ad cuius adventum omnes homines resurgere habent cum corporibus suis, et reddituri sunt de factis propriis rationem: Et qui bona egerunt, ibunt in vitam aeternam: qui vero mala, in ignem aeternum}---``At whose coming all men are to rise again with their bodies, and are to give an account of their own deeds: and those who have done good will go into eternal life, while those who have done evil will go into the everlasting fire.''\endnote{Athanasian Creed (\textit{Quicumque}), DS 76.}

In the sixth century, the Church confronted the speculative proposal associated with Origenist monks that the punishment of demons and wicked men would endure only for a limited age, ending at last in a universal restoration, or \textit{apokatastasis}. The Synod of Constantinople in 543 answered with a direct anathema in its ninth canon: ``If anyone says or holds that the punishment of the demons and of impious men is temporary, and that it will have an end at some time, that is to say, there will be a complete restoration of the demons or of impious men, let him be anathema.''\endnote{Synod of Constantinople (543), canon 9, DS 411.} Received by Pope Vigilius and confirmed by the universal Church's enduring consensus, that judgment closed the door once and for all against any doctrine that turns Gehenna into a temporary purgatory.

Seven centuries later, facing Cathar dualism---which attributed evil and the material world to a rival dark deity and denied the resurrection of the flesh---the Fourth Lateran Council (1215) promulgated its great profession of faith, \textit{Firmiter credimus}. The council confesses that the devil and the other demons were created good in nature by God and became evil by themselves, that man sinned at the devil's prompting, and that the incarnate Son of God will come at the end of the age to judge the living and the dead:

\begin{quote}
\dots\ who all shall rise with their own bodies which they now bear, that they may receive according to their works, whether they have been good or evil, the reprobate perpetual punishment with the devil, and the elect everlasting glory with Christ.\endnote{Fourth Lateran Council (1215), Constitution 1 (\textit{Firmiter credimus}), DS 801: \textit{qui omnes cum suis propriis resurgent corporibus, quae nunc gestant, ut recipiant secundum opera sua, sive bona fuerint sive mala, illi cum diabolo poenam perpetuam, et isti cum Christo gloriam sempiternam}.}
\end{quote}

Every word of that sentence carries weight. The resurrection embraces ``both reprobate and elect'' (\textit{tam reprobis quam electis}); the bodies raised are ``their own bodies which they now bear''; the retribution corresponds to their works; and the perpetuity of the reprobate's punishment with the devil stands in exact parallel with the everlasting glory of the elect with Christ.

A second doctrinal question arose during the thirteenth and fourteenth centuries: do the souls of the dead enter immediately upon their eternal reward or punishment after death, or do they wait in some suspended state until the general resurrection? The Church answered three times at the highest level of authority. First, the Second Council of Lyons (1274), receiving the profession of faith of Emperor Michael VIII Palaeologus, confessed that ``the souls of those who die in mortal sin or with original sin only descend immediately into hell, to be punished, however, with different punishments.''\endnote{Second Council of Lyons (1274), Profession of Faith of Michael Palaeologus, DS 858: \textit{Illorum autem animas, qui in mortali peccato vel cum solo originali decedunt, mox in infernum descendere, poenis tamen disparibus puniendas}.}

Sixty-two years later, after a sharp controversy sparked by Pope John XXII's personal theological sermons suggesting that the beatific vision was deferred until the Last Judgment, his successor Pope Benedict XII settled the entire question by a solemn dogmatic definition in the apostolic constitution \textit{Benedictus Deus} (29 January 1336). Having defined that the purified souls of the just behold the divine essence immediately after death, Benedict XII turned to the lost:

\begin{quote}
Moreover we define that according to the general disposition of God, the souls of those who die in actual mortal sin go down into hell immediately after death, where they are tormented by infernal punishments, and that nevertheless on the day of judgment all men will appear with their bodies ``before the tribunal of Christ'' to render an account of their own deeds, ``so that each one may receive what he has won through the body, according to what he has done, whether good or evil'' (2 Cor 5:10).\endnote{Pope Benedict XII, Apostolic Constitution \textit{Benedictus Deus} (29 January 1336), DS 1002: \textit{Diffinimus insuper, quod secundum Dei ordinationem communem animae decedentium in actuali mortali peccato mox post mortem suam ad inferna descendunt, ubi poenis infernalibus cruciantur\dots}}
\end{quote}

A century later, the Council of Florence repeated the same confession in two bulls of union with Eastern Churches. In \textit{Laetentur caeli} (6 July 1439), the Decree of Union with the Greeks, the council reaffirmed the exact wording of Lyons II: souls departing in actual mortal sin, or in original sin alone, descend immediately into hell, ``to be punished, however, with unequal punishments'' (\textit{poenis tamen disparibus puniendas}).\endnote{Council of Florence, Bull of Union with the Greeks \textit{Laetentur caeli} (6 July 1439), DS 1306.} Three years later, in \textit{Cantate Domino} (4 February 1442)---promulgated specifically as the Bull of Union with the Copts and Ethiopians, traditionally styled the Decree for the Jacobites---the council confessed the goodness of every created nature against dualism and solemnly declared that no one remaining outside the unity of the Catholic Church can become a sharer in eternal life unless joined to her before the end of life, a teaching that the Church reads together with her later clarifications on invincible ignorance and God's universal salvific will.\endnote{Council of Florence, Bull of Union with the Copts and Ethiopians \textit{Cantate Domino} (4 February 1442), DS 1333, 1351.} Notice also the careful precision of \textit{poenis tamen disparibus}: hell is not an undifferentiated leveling of every guilt. Just as Christ declared that it will be ``more tolerable'' for Tyre, Sidon, and Sodom on the day of judgment than for Corozain, Bethsaida, and Capharnaum (Matt 11:22--24), so divine justice proportions the severity of punishment to the gravity of sin.

In the sixteenth century, the Council of Trent guarded the moral and spiritual role of the fear of hell. Against Reformers who held that fear of eternal punishment turns a penitent into a hypocrite and a worse sinner, Trent declared in Session XIV, Chapter 4 and defined in Canon 5 (1551) that imperfect contrition conceived from the ugliness of sin or from the fear of hell and its punishments (\textit{ex gehennae et poenarum metu}), when joined to the hope of pardon and the exclusion of the will to sin, is a true and profitable sorrow, ``a gift of God, and an impulse of the Holy Ghost,'' preparing the sinner for justification.\endnote{Council of Trent, Session XIV (25 November 1551), Doctrine on the Sacrament of Penance, ch.\ 4, and Canon 5, DS 1678, 1705.}

The contemporary Magisterium has neither softened nor retired a syllable of this confession. Closing the Year of Faith on 30 June 1968, St.\ Paul VI proclaimed in the \textit{Credo of the People of God} (\textit{Solemni hac liturgia}) that Christ will come again in glory to judge the living and the dead, ``each according to his merits---those who have responded to the love and piety of God going to eternal life, those who have refused them to the end going to the fire that is not extinguished.''\endnote{Pope St.\ Paul VI, \textit{Solemni hac liturgia} (30 June 1968), \S 12.} The \textit{Catechism of the Catholic Church} sums up the entire inheritance in five dense paragraphs (CCC 1033--1037) and in its treatment of mortal sin (CCC 1861):

\begin{quote}
The teaching of the Church affirms the existence of hell and its eternity. Immediately after death the souls of those who die in a state of mortal sin descend into hell, where they suffer the punishments of hell, ``eternal fire.'' The chief punishment of hell is eternal separation from God, in whom alone man can possess the life and happiness for which he was created and for which he longs. (CCC 1035)\endnote{\textit{Catechism of the Catholic Church}, \S\S 1033--1037, 1861.}
\end{quote}

To this the Catechism joins two truths that protect the heart from fatalism: ``God predestines no one to go to hell; for this, a willful turning away from God (a mortal sin) is necessary, and persistence in it to the end'' (CCC 1037). St.\ John Paul II pressed the same point home in his General Audience catechesis of 28 July 1999: hell is not an arbitrary penalty inflicted from outside by an angry deity, but ``a development of premises already set by people in this life''---``the state of those who freely and definitively separate themselves from God, the source of all life and joy.''\endnote{Pope St.\ John Paul II, General Audience, 28 July 1999.}

Taken together, the councils and the Catechism outline the fixed shape of Catholic belief: hell exists; its punishment is everlasting; those who die in unrepented mortal sin enter it immediately upon death; their punishments are unequal according to guilt; at the Last Day they will rise in their own bodies to stand before Christ; God predestines no one to this ruin; and the dread of it, far from corrupting the conscience, can be a gift of the Holy Ghost drawing the soul back to grace.
